\documentclass[12pt]{article}
\usepackage[utf8]{inputenc}
\usepackage[T1]{fontenc}
\usepackage{hyperref}
\title{Shumashti Language Resources for K-12 Curriculum Development}
\author{Compiled by Claude Sonnet 5.5 (claude-sonnet-5-5) for Andr\'as Kornai}
\date{}
\begin{document}
\maketitle

\section{Introduction}
Shumashti (Shumasht; ISO 639-3 sts; Glottocode shum1235) is a Dardic Indo-Aryan language of Kunar Province in eastern Afghanistan. A few speakers are reported in Chitral, Pakistan. Estimates range from about 1,000 to 1,600 speakers. UNESCO classes it as severely endangered, and one source says it is spoken in only a few households. It is unwritten. The searches found no educational materials of any kind.

\section{Primary Educational Resources}
None found. No primers, textbooks or school programmes exist in the search results.

\section{Linguistic and Reference Materials}
Only general descriptions were found.
\begin{itemize}
\item Overview of the language (Wikipedia). \url{https://en.wikipedia.org/wiki/Shumashti_language}
\item Endangered Languages Project record. \url{https://endangeredlanguages.com/lang/5425}
\item Verbix endangered-languages page. \url{https://docs.verbix.com/EndangeredLanguages/Shumashti}
\item Rosetta Project (Long Now) entry. \url{https://rosettapanglossia.longnow.org/index.php?title=Sts}
\end{itemize}

\section{Religious and Community Resources}
None found.

\section{Digital and Audio Resources}
No dictionaries, corpora or recordings were found.

\section{Curriculum Development Context}
The language has no writing system and no found teaching materials. A curriculum would require fieldwork, an orthography decided with the community, and basic documentation first.

\end{document}